\begin{recipe}
[% 
    preparationtime = {\unit[15]{min}},
    portion = {\portion{1}},
    source = {Youtube: Kwoowk, das exakte Video finde ich aber nicht mehr},
    calory = {510 kcal | 26g Protein | 49g KH | 25g Fett},
    rating = {1},
    price = {1,75€}
]
{Savory Oatmeal}

% \graph
% {% pictures
%     small=pic/oatmeal_small,
%     big=pic/oatmeal_big
% }

\introduction{%
    Ein Frühstück muss nicht immer süß sein, hier ein herzhaftes Oatmeal, welches den Magen schnell füllt. Ich muss zwar sagen, dass es sehr gewöhnungsbedürftig ist, aber man kann es ja mal ausprobieren.
}

\ingredients{%
    70 g & Haferflocken\index{Getreide, Pasta \& Brot!Haferflocken}\\
    290 ml & Wasser\index{Getränke!Wasser}\\
    1 EL & Sojasauce\index{Öle, Saucen \& Würzmittel!Sojasauce}\\
    1 TL & Miso-Paste\index{Öle, Saucen \& Würzmittel!Miso}\\
    2 Stangen & Frühlingszwiebeln\index{Gemüse \& Pilze!Frühlingszwiebeln}\\
    \nicefrac{1}{2} EL & Sesamöl\index{Öle, Saucen \& Würzmittel!Sesamöl}\\
    2 & Eier (optional)\index{Eier!Eier}\\
    1 TL & Sesam (optional)\index{Nüsse \& Samen!Sesam}
}

\preparation{%
    \step%
    Haferflocken in einem Topf ohne Öl kurz anrösten, bis sie leicht gebräunt sind.
    
    \step%
    Wasser, Sojasauce und Miso-Paste dazugeben und gut umrühren.
    
    \step%
    Köcheln lassen, bis die gewünschte Konsistenz erreicht ist (je nach Haferflocken ein paar Minuten).
    
    \step%
    Frühlingszwiebeln hacken, mit Sesamöl in einer Pfanne kurz (ca. 1 Minute) anbraten und unterrühren oder oben drauf geben.
    
    \step%
    In eine Schüssel geben und servieren.\\
    
    \step%
    Optional: Spiegelei braten, auf das Oatmeal legen und mit Sesam toppen.
}

\suggestion{%
    Wenn man es schärfer will: Chiliöl oder Sriracha passt extrem gut.
}

\hint{%
    Miso nicht zu lange hart kochen lassen, lieber am Ende einrühren, dann bleibt der Geschmack runder.
}
    
\end{recipe}